% Integration fragment: finite feedback, not Lean formalization.
\subsection{Four certified finite feedback steps}
\label{sec:second-seed-certified-finite-feedback}

The positive operator of Section~\ref{sec:second-seed-finite-feedback}
can be used a fixed number of times with an exactly certified input.
The construction below uses the original coefficient $K=681/100$
throughout. It does not depend on a numerical infinite iteration or
on the applicability of the newer modulus weights at a smaller cutoff.

Write $S=11/5$, $s_0=227/50$, and
$(\phi_1,\phi_2,\phi_3)=(353/100,29/10,61/25)$.
The family consists of the typed stable-prefix and classical states
of Section~\ref{sec:new-finite-double-sieve-recursion}; in particular
the seed comparison cutoff must satisfy $L/S<\chi$ on enhanced
states. Every descendant retains the full and selected square tests,
the actual residual level, its stable ceiling, and $\alpha\ge2$.
No assertion below discards this comparison-cutoff condition.

\begin{proposition}[A larger seed after four finite rounds]
\label{prop:second-seed-certified-finite-feedback}
Under the analytic inputs of
Proposition~\ref{prop:second-seed-stable-prefix}, on the admitted
states just specified, one has
\[
 \Phi_{11/5}\le
 \left(1-\frac{2327}{100000}+o(1)\right)\Theta_Q.
\]
The estimate is uniform in the admitted $\alpha\ge2$ and uses only
four finite feedback rounds starting from the seed $1/50$.
\end{proposition}

\begin{proof}
For a previously proved positive constant $J$, choose only the
following upper defect and its one-step Buchstab lower consequence:
\[
 H_J(w)=J\mathbf1_{[1,S]}(w),\qquad
 h_J(w)=J\log\frac{S}{w-1}\,\mathbf1_{[2,1+S]}(w).
\]
The upper option on $w<S$ follows by comparing positive counts
with the admitted smaller cutoff at $S$; the ordinary normalized
upper multiplier is one on this range. The lower option is
$\int_{w-1}^S J\,dt/t$, including the already justified upper
children with $1\le t<2$. It is used only on states satisfying
their inherited comparison and square conditions.

Here $h_J(w)$ needs admission of its own physical cutoff, not a
new assumption that this lower parent itself admits argument $S$.
Indeed, for a Buchstab child with argument $v\in[w-1,S]$, its new
prime has exponent $L/(v+1)$ and its residual exponent is
$Lv/(v+1)$. Its seed comparison exponent satisfies
\[
 \frac{Lv}{S(v+1)}\le\frac{L}{v+1},
\]
so the compared cutoff is at most that new prime, and remains
below the old stable ceiling. Since $v\ge1$, its new square test
also holds. This distinction is needed for the single-prime lower
children of the weighted inequality: they need not themselves
admit the larger cutoff associated with argument $S$.

Insert these previously proved options in
\eqref{eq:second-seed-feedback-operator}. The $\Gamma_1$ feedback
vanishes, since $s_0,\phi_1>S$. The remaining feedback is a positive
linear function of $J$; write it as $\lambda J$.
The three single-prime integrals use
$\rho=s_0(1-u)$ and hence start at
\[
 u_* = 1-\frac{1+S}{s_0}=\frac{67}{227}.
\]
Their upper endpoints are $1/S,1/\phi_2,1/\phi_3$ respectively.
All their actual arguments exceed two, so no lower defect below
two is used.

For the four pair terms, set $a_i=1/\phi_i$ and $a_0=1/s_0$.
The exact support restrictions for the constant upper defect are
\[
 u+v\ge\frac{117}{227}\quad(i=5,6),\qquad
 v\ge1-\frac{16}{5}u\quad(i=7,8).
\]
They are intersected with the original $D_5^+,D_6,D_7,D_8$;
there is no feedback on the bad part $D_5^-$.
For fixed $u$, if the resulting lower and upper endpoints are
$l(u),r(u)$, the exact pair contribution is
\[
 J\int \frac{1}{u(1-u)}
 \log\frac{r(u)(1-u-l(u))}{l(u)(1-u-r(u))}\,du.
\]
All endpoints are affine with rational coefficients. For $D_5^+$
the three $u$ intervals end at $117/454,19/58,1/3$, with initial
endpoint $50/227$. The rectangle $D_6$ needs no support truncation.
For $D_7$ the nonempty support starts at $1265/5648$ and its lower
boundary changes at $5/21$. For $D_8$ the lower boundary changes
at $1265/5648$. These cuts also follow directly by intersecting
the displayed lines; no floating boundary decision is required.

The certificate \texttt{certify\_second\_seed\_first\_feedback.py}
uses 4096 rational Darboux intervals on each resulting piece and
the previously audited 80-bit outward logarithm arithmetic. At
$J=1/50$, after division by five, it gives
\[
 \lambda\frac1{50}\ge
 \frac{3152244426576931077198}{1208925819614629174706176}
 >\frac{13}{5000},
 \qquad\text{hence}\qquad \lambda>\frac{13}{100}.
\]
Separately, the ordinary seed certificate
\eqref{eq:second-seed-omega-exact-bound} gives
\[
 J_2^\omega\ge
 \frac{24479824718463837928127}{1208925819614629174706176}
 >\frac{20249}{10^6}.
\]
Consequently an admitted uniform short-support seed $J_n$ proves
the next uniform seed
\[
 J_{n+1}=\frac{20249}{10^6}+\frac{13}{100}J_n.
\]
At each round first derive $h_{J_n}$ from the already proved
$H_{J_n}$, and then apply the second-seed inequality. Thus no
unknown estimate from the new round is used to prove itself.
Positive-count comparison then supplies $H_{J_{n+1}}$ on the
short interval, with the same comparison-cutoff qualification.

Starting with $J_0=1/50$, the four exact rational values are
\[
 \frac{22849}{10^6},\quad
 \frac{2321937}{10^8},\quad
 \frac{232675181}{10^{10}},\quad
 \frac{23273777353}{10^{12}}
 >\frac{2327}{100000}.
\]
They are checked in
\texttt{certify\_second\_seed\_finite\_iteration.json}.
There are only finitely many Buchstab and weighted steps. On
compact pieces their new prime exponents remain bounded below by
a fixed positive constant depending on this finite depth. The
already proved fixed-positive-cutoff product theorem with $K=6.81$
therefore applies after its precision parameters are chosen for
that depth. Finitely many errors remain $o(1)$, and the strict
final slack absorbs them in the inherited order of limits.
\end{proof}

For insertion in the earlier narrow lower-integral block, first
use the existing stable descendant propagation at argument three.
Retain the original threshold condition $h>3/20$, so the entire
anchor interval $29/200<t<3/20$ remains inside that lower integral.
It is linear in the short-support seed and gives
\[
 H(3)>\frac{28}{333}\frac{2327}{100000}.
\]
The second narrow-block contribution thus becomes
\[
 \frac{112}{837495}\,
 \frac{2327/100000}{1/50}.
\]
This replaces the older value for the same block. It is not added
to it. In particular, the seed at argument $11/5$ is not applied
directly to an anchor whose comparison cutoff exceeds its smallest
specified prime; the separate argument-three propagation is
essential. No entire floating $H,h$ table or global Chen coefficient
is asserted by this finite seed certificate alone.
